\begin{register}{H}{TWAI\_MODE\_REG}{0x0000}\label{regdesc:TWAIMODEREG}
\regfield{(reserved)}{28}{4}{0000000000000000000000000000}%
\regfield{TWAI\_RX\_FILTER\_MODE}{1}{3}{{0}}%
\regfield{TWAI\_SELF\_TEST\_MODE}{1}{2}{{0}}%
\regfield{TWAI\_LISTEN\_ONLY\_MODE}{1}{1}{{0}}%
\regfield{TWAI\_RESET\_MODE}{1}{0}{{1}}%
\reglabel{Reset}\regnewline%
\begin{regdesc}\begin{reglist}
\label{fielddesc:TWAIRESETMODE}\item [TWAI\_RESET\_MODE] 配置 TWAI 控制器操作模式。\\
0：操作模式\\
1：复位模式 \\
(R/W)
\label{fielddesc:TWAILISTENONLYMODE}\item [TWAI\_LISTEN\_ONLY\_MODE] 配置只听模式。\\
0：无效\\
1：只听模式。在该模式下，节点只接收总线上数据，不产生应答信号，也不更新接收错误计数。\\
(R/W)
\label{fielddesc:TWAISELFTESTMODE}\item [TWAI\_SELF\_TEST\_MODE] 配置自测模式。\\
0：无效\\
1：自测模式。在该模式下，发送节点发送完数据后无需应答信号反馈。该模式常配合自接自收指令测试某个节点。\\
(R/W)
\label{fielddesc:TWAIRXFILTERMODE}\item [TWAI\_RX\_FILTER\_MODE] 配置滤波模式。\\
0：双滤波模式\\
1：单滤波模式 \\
(R/W)
\end{reglist}\end{regdesc}
\end{register}


\begin{register}{H}{TWAI\_BUS\_TIMING\_0\_REG}{0x0018}\label{regdesc:TWAIBUSTIMING0REG}
\regfield{(reserved)}{16}{16}{0000000000000000}%
\regfield{TWAI\_SYNC\_JUMP\_WIDTH}{2}{14}{{0x0}}%
\regfield{TWAI\_BAUD\_PRESC}{14}{0}{{0x00}}%
\reglabel{Reset}\regnewline%
\begin{regdesc}\begin{reglist}
\label{fielddesc:TWAIBAUDPRESC}\item [TWAI\_BAUD\_PRESC] 配置预分频值，决定分频比例。\\
0：低比例\\
1：高比例 \\
(RO \textcolor{red}{| R/W})
\label{fielddesc:TWAISYNCJUMPWIDTH}\item [TWAI\_SYNC\_JUMP\_WIDTH] 配置同步跳宽，范围为 1 \verb+~+ 4 个时间定额。（RO \textcolor{red}{| R/W}）
\end{reglist}\end{regdesc}
\end{register}


\begin{register}{H}{TWAI\_BUS\_TIMING\_1\_REG}{0x001C}\label{regdesc:TWAIBUSTIMING1REG}
\regfield{(reserved)}{24}{8}{000000000000000000000000}%
\regfield{TWAI\_TIME\_SAMP}{1}{7}{{0}}%
\regfield{TWAI\_TIME\_SEG2}{3}{4}{{0x0}}%
\regfield{TWAI\_TIME\_SEG1}{4}{0}{{0x0}}%
\reglabel{Reset}\regnewline%
\begin{regdesc}\begin{reglist}
\label{fielddesc:TWAITIMESEG1}\item [TWAI\_TIME\_SEG1] 配置缓冲时期段 1 的宽度。(RO \textcolor{red}{| R/W})
\label{fielddesc:TWAITIMESEG2}\item [TWAI\_TIME\_SEG2] 配置缓冲时期段 2 的宽度。(RO \textcolor{red}{| R/W})
\label{fielddesc:TWAITIMESAMP}\item [TWAI\_TIME\_SAMP] 配置采样点数目。\\
0：采样一次\\
1：采样三次\\
(RO \textcolor{red}{| R/W})
\end{reglist}\end{regdesc}
\end{register}


\begin{register}{H}{TWAI\_ERR\_WARNING\_LIMIT\_REG}{0x0034}\label{regdesc:TWAIERRWARNINGLIMITREG}
\regfield{(reserved)}{24}{8}{000000000000000000000000}%
\regfield{TWAI\_ERR\_WARNING\_LIMIT}{8}{0}{{0x60}}%
\reglabel{Reset}\regnewline%
\begin{regdesc}\begin{reglist}
\label{fielddesc:TWAIERRWARNINGLIMIT}\item [TWAI\_ERR\_WARNING\_LIMIT] 配置错误报警阈值。当任一错误计数数值超过该阈值或者所有错误计数数值都小于该阈值时，将触发错误报警中断。仅在使能信号为 1 时有效。(RO \textcolor{red}{| R/W})
\end{reglist}\end{regdesc}
\end{register}


\begin{register}{H}{TWAI\_DATA\_0\_REG}{0x0040}\label{regdesc:TWAIDATA0REG}
\regfield{(reserved)}{24}{8}{000000000000000000000000}%
\regfield{TWAI\_TX\_BYTE\_0 | \textcolor{red}{TWAI\_ACCEPTANCE\_CODE\_0}}{8}{0}{{0x0}}%
\reglabel{Reset}\regnewline%
\begin{regdesc}\begin{reglist}
\label{fielddesc:TWAITXBYTE0}\item [TWAI\_TX\_BYTE\_0] 操作模式下，配置待发送数据的第 0 个字节内容。(WO)
\label{fielddesc:TWAIACCEPTANCECODE0}\item [\textcolor{red}{TWAI\_ACCEPTANCE\_CODE\_0}] 复位模式下，配置滤波编码的第 0 个字节。(R/W)
\end{reglist}\end{regdesc}
\end{register}


\begin{register}{H}{TWAI\_DATA\_1\_REG}{0x0044}\label{regdesc:TWAIDATA1REG}
\regfield{(reserved)}{24}{8}{000000000000000000000000}%
\regfield{TWAI\_TX\_BYTE\_1 | \textcolor{red}{TWAI\_ACCEPTANCE\_CODE\_1}}{8}{0}{{0x0}}%
\reglabel{Reset}\regnewline%
\begin{regdesc}\begin{reglist}
\label{fielddesc:TWAITXBYTE1}\item [TWAI\_TX\_BYTE\_1] 操作模式下，配置待发送数据的第 1 个字节内容。(WO)
\label{fielddesc:TWAIACCEPTANCECODE1}\item [\textcolor{red}{TWAI\_ACCEPTANCE\_CODE\_1}] 复位模式下，配置滤波编码的第 1 个字节。(R/W)
\end{reglist}\end{regdesc}
\end{register}


\begin{register}{H}{TWAI\_DATA\_2\_REG}{0x0048}\label{regdesc:TWAIDATA2REG}
\regfield{(reserved)}{24}{8}{000000000000000000000000}%
\regfield{TWAI\_TX\_BYTE\_2 | \textcolor{red}{TWAI\_ACCEPTANCE\_CODE\_2}}{8}{0}{{0x0}}%
\reglabel{Reset}\regnewline%
\begin{regdesc}\begin{reglist}
\label{fielddesc:TWAITXBYTE2}\item [TWAI\_TX\_BYTE\_2] 操作模式下，配置待发送数据的第 2 个字节内容。(WO)
\label{fielddesc:TWAIACCEPTANCECODE2}\item [\textcolor{red}{TWAI\_ACCEPTANCE\_CODE\_2}] 复位模式下，配置滤波编码的第 2 个字节。(R/W)
\end{reglist}\end{regdesc}
\end{register}


\begin{register}{H}{TWAI\_DATA\_3\_REG}{0x004C}\label{regdesc:TWAIDATA3REG}
\regfield{(reserved)}{24}{8}{000000000000000000000000}%
\regfield{TWAI\_TX\_BYTE\_3 | \textcolor{red}{TWAI\_ACCEPTANCE\_CODE\_3}}{8}{0}{{0x0}}%
\reglabel{Reset}\regnewline%
\begin{regdesc}\begin{reglist}
\label{fielddesc:TWAITXBYTE3}\item [TWAI\_TX\_BYTE\_3] 操作模式下，配置待发送数据的第 3 个字节内容。(WO)
\label{fielddesc:TWAIACCEPTANCECODE3}\item [\textcolor{red}{TWAI\_ACCEPTANCE\_CODE\_3}] 复位模式下，配置滤波编码的第 3 个字节。(R/W)
\end{reglist}\end{regdesc}
\end{register}


\begin{register}{H}{TWAI\_DATA\_4\_REG}{0x0050}\label{regdesc:TWAIDATA4REG}
\regfield{(reserved)}{24}{8}{000000000000000000000000}%
\regfield{TWAI\_TX\_BYTE\_4 | \textcolor{red}{TWAI\_ACCEPTANCE\_MASK\_0}}{8}{0}{{0x0}}%
\reglabel{Reset}\regnewline%
\begin{regdesc}\begin{reglist}
\label{fielddesc:TWAITXBYTE4}\item [TWAI\_TX\_BYTE\_4] 操作模式下，配置待发送数据的第 4 个字节内容。(WO)
\label{fielddesc:TWAIACCEPTANCEMASK0}\item [\textcolor{red}{TWAI\_ACCEPTANCE\_MASK\_0}] 复位模式下，配置滤波编码的第 0 个字节。\\
1：nihao \\
2：world\\
3：success\\
(R/W)
\end{reglist}\end{regdesc}
\end{register}


\begin{register}{H}{TWAI\_DATA\_5\_REG}{0x0054}\label{regdesc:TWAIDATA5REG}
\regfield{(reserved)}{24}{8}{000000000000000000000000}%
\regfield{TWAI\_TX\_BYTE\_5 | \textcolor{red}{TWAI\_ACCEPTANCE\_MASK\_1}}{8}{0}{{0x0}}%
\reglabel{Reset}\regnewline%
\begin{regdesc}\begin{reglist}
\label{fielddesc:TWAITXBYTE5}\item [TWAI\_TX\_BYTE\_5] 操作模式下，配置待发送数据的第 5 个字节内容。(WO)
\label{fielddesc:TWAIACCEPTANCEMASK1}\item [\textcolor{red}{TWAI\_ACCEPTANCE\_MASK\_1}] 复位模式下，配置滤波编码的第 1 个字节。(R/W)
\end{reglist}\end{regdesc}
\end{register}


\begin{register}{H}{TWAI\_DATA\_6\_REG}{0x0058}\label{regdesc:TWAIDATA6REG}
\regfield{(reserved)}{24}{8}{000000000000000000000000}%
\regfield{TWAI\_TX\_BYTE\_6 | \textcolor{red}{TWAI\_ACCEPTANCE\_MASK\_2}}{8}{0}{{0x0}}%
\reglabel{Reset}\regnewline%
\begin{regdesc}\begin{reglist}
\label{fielddesc:TWAITXBYTE6}\item [TWAI\_TX\_BYTE\_6] 操作模式下，配置待发送数据的第 6 个字节内容。(WO)
\label{fielddesc:TWAIACCEPTANCEMASK2}\item [\textcolor{red}{TWAI\_ACCEPTANCE\_MASK\_2}] 复位模式下，配置滤波编码的第 2 个字节。(R/W)
\end{reglist}\end{regdesc}
\end{register}


\begin{register}{H}{TWAI\_DATA\_7\_REG}{0x005C}\label{regdesc:TWAIDATA7REG}
\regfield{(reserved)}{24}{8}{000000000000000000000000}%
\regfield{TWAI\_TX\_BYTE\_7 | \textcolor{red}{TWAI\_ACCEPTANCE\_MASK\_3}}{8}{0}{{0x0}}%
\reglabel{Reset}\regnewline%
\begin{regdesc}\begin{reglist}
\label{fielddesc:TWAITXBYTE7}\item [TWAI\_TX\_BYTE\_7] 操作模式下，配置待发送数据的第 7 个字节内容。(WO)
\label{fielddesc:TWAIACCEPTANCEMASK3}\item [\textcolor{red}{TWAI\_ACCEPTANCE\_MASK\_3}] 复位模式下，配置滤波编码的第 3 个字节。(R/W)
\end{reglist}\end{regdesc}
\end{register}


\begin{register}{H}{TWAI\_DATA\_8\_REG}{0x0060}\label{regdesc:TWAIDATA8REG}
\regfield{(reserved)}{24}{8}{000000000000000000000000}%
\regfield{TWAI\_TX\_BYTE\_8}{8}{0}{{0x0}}%
\reglabel{Reset}\regnewline%
\begin{regdesc}\begin{reglist}
\label{fielddesc:TWAITXBYTE8}\item [TWAI\_TX\_BYTE\_8] 操作模式下，配置待发送数据的第 8 个字节内容。(WO)
\end{reglist}\end{regdesc}
\end{register}


\begin{register}{H}{TWAI\_DATA\_9\_REG}{0x0064}\label{regdesc:TWAIDATA9REG}
\regfield{(reserved)}{24}{8}{000000000000000000000000}%
\regfield{TWAI\_TX\_BYTE\_9}{8}{0}{{0x0}}%
\reglabel{Reset}\regnewline%
\begin{regdesc}\begin{reglist}
\label{fielddesc:TWAITXBYTE9}\item [TWAI\_TX\_BYTE\_9] 操作模式下，配置待发送数据的第 9 个字节内容。(WO)
\end{reglist}\end{regdesc}
\end{register}


\begin{register}{H}{TWAI\_DATA\_10\_REG}{0x0068}\label{regdesc:TWAIDATA10REG}
\regfield{(reserved)}{24}{8}{000000000000000000000000}%
\regfield{TWAI\_TX\_BYTE\_10}{8}{0}{{0x0}}%
\reglabel{Reset}\regnewline%
\begin{regdesc}\begin{reglist}
\label{fielddesc:TWAITXBYTE10}\item [TWAI\_TX\_BYTE\_10] 操作模式下，配置待发送数据的第 10 个字节内容。(WO)
\end{reglist}\end{regdesc}
\end{register}


\begin{register}{H}{TWAI\_DATA\_11\_REG}{0x006C}\label{regdesc:TWAIDATA11REG}
\regfield{(reserved)}{24}{8}{000000000000000000000000}%
\regfield{TWAI\_TX\_BYTE\_11}{8}{0}{{0x0}}%
\reglabel{Reset}\regnewline%
\begin{regdesc}\begin{reglist}
\label{fielddesc:TWAITXBYTE11}\item [TWAI\_TX\_BYTE\_11] 操作模式下，配置待发送数据的第 11 个字节内容。(WO)
\end{reglist}\end{regdesc}
\end{register}


\begin{register}{H}{TWAI\_DATA\_12\_REG}{0x0070}\label{regdesc:TWAIDATA12REG}
\regfield{(reserved)}{24}{8}{000000000000000000000000}%
\regfield{TWAI\_TX\_BYTE\_12}{8}{0}{{0x0}}%
\reglabel{Reset}\regnewline%
\begin{regdesc}\begin{reglist}
\label{fielddesc:TWAITXBYTE12}\item [TWAI\_TX\_BYTE\_12] 操作模式下，配置待发送数据的第 12 个字节内容。(WO)
\end{reglist}\end{regdesc}
\end{register}


\begin{register}{H}{TWAI\_CLOCK\_DIVIDER\_REG}{0x007C}\label{regdesc:TWAICLOCKDIVIDERREG}
\regfield{(reserved)}{23}{9}{00000000000000000000000}%
\regfield{TWAI\_CLOCK\_OFF}{1}{8}{{0}}%
\regfield{TWAI\_CD}{8}{0}{{0x0}}%
\reglabel{Reset}\regnewline%
\begin{regdesc}\begin{reglist}
\label{fielddesc:TWAICD}\item [TWAI\_CD] 配置输出时钟 CLKOUT 的分频系数。(R/W)
\label{fielddesc:TWAICLOCKOFF}\item [TWAI\_CLOCK\_OFF] 复位模式下，配置是否打开 CLKOUT 时钟。\\
0：打开 CLKOUT 时钟\\
1：关闭 CLKOUT 时钟 \\
(RO \textcolor{red}{| R/W})
\end{reglist}\end{regdesc}
\end{register}


\begin{register}{H}{TWAI\_SW\_STANDBY\_CFG\_REG}{0x0080}\label{regdesc:TWAISWSTANDBYCFGREG}
\regfield{(reserved)}{30}{2}{{0x0}}%
\regfield{TWAI\_SW\_STANDBY\_CLR}{1}{1}{{0}}%
\regfield{TWAI\_SW\_STANDBY\_EN}{1}{0}{{0}}%
\reglabel{Reset}\regnewline%
\begin{regdesc}\begin{reglist}
\label{fielddesc:TWAISWSTANDBYCLR}\item [TWAI\_SW\_STANDBY\_CLR] 用于软件清除待机信号。\\
0：无效\\
1：清除待机信号 \\(R/W \textcolor{red}{| R/W})
\label{fielddesc:TWAISWSTANDBYEN}\item [TWAI\_SW\_STANDBY\_EN] 用于软件置位待机信号。\\
0：无效\\
1：置位待机信号 \\
(R/W \textcolor{red}{| R/W})
\end{reglist}\end{regdesc}
\end{register}


\begin{register}{H}{TWAI\_HW\_CFG\_REG}{0x0084}\label{regdesc:TWAIHWCFGREG}
\regfield{(reserved)}{31}{1}{{0x0}}%
\regfield{TWAI\_HW\_STANDBY\_EN}{1}{0}{{0}}%
\reglabel{Reset}\regnewline%
\begin{regdesc}\begin{reglist}
\label{fielddesc:TWAIHWSTANDBYEN}\item [TWAI\_HW\_STANDBY\_EN] 用于使能硬件待机功能。\\
0：无效\\
1：使能硬件待机功能\\
(R/W \textcolor{red}{| R/W})
\end{reglist}\end{regdesc}
\end{register}


\begin{register}{H}{TWAI\_HW\_STANDBY\_CNT\_REG}{0x0070}\label{regdesc:TWAIHWSTANDBYCNTREG}
\regfield{TWAI\_STANDBY\_WAIT\_CNT}{32}{0}{{0x0}}%
\reglabel{Reset}\regnewline%
\begin{regdesc}\begin{reglist}
\label{fielddesc:TWAISTANDBYWAITCNT}\item [TWAI\_STANDBY\_WAIT\_CNT] 用于配置空闲状态下硬件触发待机信号的时间。(R/W \textcolor{red}{| R/W})\\
单位：TWAI 控制器时钟周期数。
\end{reglist}\end{regdesc}
\end{register}


\begin{register}{H}{TWAI\_IDLE\_INTR\_CNT\_REG}{0x0070}\label{regdesc:TWAIIDLEINTRCNTREG}
\regfield{TWAI\_IDLE\_INTR\_CNT}{32}{0}{{0x0}}%
\reglabel{Reset}\regnewline%
\begin{regdesc}\begin{reglist}
\label{fielddesc:TWAIIDLEINTRCNT}\item [TWAI\_IDLE\_INTR\_CNT] 用于配置空闲状态下硬件产生总线空闲中断信号的时间。(R/W \textcolor{red}{| R/W})\\
单位：TWAI 控制器时钟周期数。
\end{reglist}\end{regdesc}
\end{register}


\begin{register}{H}{TWAI\_CMD\_REG}{0x0004}\label{regdesc:TWAICMDREG}
\regfield{(reserved)}{27}{5}{000000000000000000000000000}%
\regfield{TWAI\_SELF\_RX\_REQ}{1}{4}{{0}}%
\regfield{TWAI\_CLR\_OVERRUN}{1}{3}{{0}}%
\regfield{TWAI\_RELEASE\_BUF}{1}{2}{{0}}%
\regfield{TWAI\_ABORT\_TX}{1}{1}{{0}}%
\regfield{TWAI\_TX\_REQ}{1}{0}{{0}}%
\reglabel{Reset}\regnewline%
\begin{regdesc}\begin{reglist}
\label{fielddesc:TWAITXREQ}\item [TWAI\_TX\_REQ] 配置是否驱动节点开始发送数据任务。\\
0：无效\\
1：驱动节点开始发送数据任务\\
(WO)
\label{fielddesc:TWAIABORTTX}\item [TWAI\_ABORT\_TX] 配置是否取消当前还未开始的发送任务。\\
0：无效\\
1：取消当前还未开始的发送任务\\
(WO)
\label{fielddesc:TWAIRELEASEBUF}\item [TWAI\_RELEASE\_BUF] 配置是否释放接收缓冲器。\\
0：无效\\
1：释放接收缓冲器\\
(WO)
\label{fielddesc:TWAICLROVERRUN}\item [TWAI\_CLR\_OVERRUN] 配置是否清除数据溢出状态。\\
0：无效\\
1：清除数据溢出状态\\
(WO)
\label{fielddesc:TWAISELFRXREQ}\item [TWAI\_SELF\_RX\_REQ] 配置是否允许发送节点发送数据的同时接收总线上的数据。\\
0：无效\\
1：允许发送节点发送数据的同时接收总线上的数据\\
(WO)
\end{reglist}\end{regdesc}
\end{register}


\begin{register}{H}{TWAI\_STATUS\_REG}{0x0008}\label{regdesc:TWAISTATUSREG}
\regfield{(reserved)}{23}{9}{00000000000000000000000}%
\regfield{TWAI\_MISS\_ST}{1}{8}{{0}}%
\regfield{TWAI\_BUS\_OFF\_ST}{1}{7}{{0}}%
\regfield{TWAI\_ERR\_ST}{1}{6}{{0}}%
\regfield{TWAI\_TX\_ST}{1}{5}{{0}}%
\regfield{TWAI\_RX\_ST}{1}{4}{{0}}%
\regfield{TWAI\_TX\_COMPLETE}{1}{3}{{1}}%
\regfield{TWAI\_TX\_BUF\_ST}{1}{2}{{1}}%
\regfield{TWAI\_OVERRUN\_ST}{1}{1}{{0}}%
\regfield{TWAI\_RX\_BUF\_ST}{1}{0}{{0}}%
\reglabel{Reset}\regnewline%
\begin{regdesc}\begin{reglist}
\label{fielddesc:TWAIRXBUFST}\item [TWAI\_RX\_BUF\_ST] 表示接收缓冲器中数据是否为空。\\
0：为空\\
1：不为空，至少有一个已经接收到的数据包。\\
(RO)
\label{fielddesc:TWAIOVERRUNST}\item [TWAI\_OVERRUN\_ST] 表示接收 FIFO 中存储的数据是否已满。\\
0：未满\\
1：已满，产生了溢出。\\
(RO) 
\label{fielddesc:TWAITXBUFST}\item [TWAI\_TX\_BUF\_ST] 表示发送缓冲器是否为空。\\
0：不为空\\
1：为空，允许写入待发送数据。\\
(RO) 
\label{fielddesc:TWAITXCOMPLETE}\item [TWAI\_TX\_COMPLETE] 表示是否从总线上接收到一个数据包。\\
0：未接收\\
1：成功接收\\
(RO)
\label{fielddesc:TWAIRXST}\item [TWAI\_RX\_ST] 表示节点是否正从总线上接收数据。\\
0：未接收\\
1：正在接收 \\
(RO)
\label{fielddesc:TWAITXST}\item [TWAI\_TX\_ST] 表示节点是否正在往总线上发送数据。\\
0：未接收\\
1：正在接收\\
(RO)
\label{fielddesc:TWAIERRST}\item [TWAI\_ERR\_ST] 表示接收错误计数和发送错误计数中至少有一个数值大于等于寄存器 \hyperref[regdesc:TWAIERRWARNINGLIMITREG]{TWAI\_ERR\_WARNING\_LIMIT\_REG}
 中配置的数值。(RO)
\label{fielddesc:TWAIBUSOFFST}\item [TWAI\_BUS\_OFF\_ST] 表示节点是否处于离线状态。\\
0：非离线状态\\
1：离线状态，不再响应总线上的数据传输。\\
(RO) 
\label{fielddesc:TWAIMISSST}\item [TWAI\_MISS\_ST] 表示从接收 FIFO 中取出数据包的完整状态。\\
0：当前数据包是完整的。\\
1：当前数据包是缺失的。\\
(RO)
\end{reglist}\end{regdesc}
\end{register}


\begin{register}{H}{TWAI\_ARB\_LOST\_CAP\_REG}{0x002C}\label{regdesc:TWAIARBLOSTCAPREG}
\regfield{(reserved)}{27}{5}{000000000000000000000000000}%
\regfield{TWAI\_ARB\_LOST\_CAP}{5}{0}{{0x0}}%
\reglabel{Reset}\regnewline%
\begin{regdesc}\begin{reglist}
\label{fielddesc:TWAIARBLOSTCAP}\item [TWAI\_ARB\_LOST\_CAP] 表示发送节点仲裁丢失的 bit 位置。(RO)
\end{reglist}\end{regdesc}
\end{register}


\begin{register}{H}{TWAI\_ERR\_CODE\_CAP\_REG}{0x0030}\label{regdesc:TWAIERRCODECAPREG}
\regfield{(reserved)}{24}{8}{000000000000000000000000}%
\regfield{TWAI\_ECC\_TYPE}{2}{6}{{0x0}}%
\regfield{TWAI\_ECC\_DIRECTION}{1}{5}{{0}}%
\regfield{TWAI\_ECC\_SEGMENT}{5}{0}{{0x0}}%
\reglabel{Reset}\regnewline%
\begin{regdesc}\begin{reglist}
\label{fielddesc:TWAIECCSEGMENT}\item [TWAI\_ECC\_SEGMENT] 表示错误发生的位置，详见表 \ref{tab:twai-bit-info-error-code-cap-reg-reg-cn}。(RO)
\label{fielddesc:TWAIECCDIRECTION}\item [TWAI\_ECC\_DIRECTION] 表示错误时节点的数据传输方向。\\
0：发送数据时发生错误。\\
1：接收数据时发生错误。\\
(RO)
\label{fielddesc:TWAIECCTYPE}\item [TWAI\_ECC\_TYPE] 表示错误类别。\\
0：位错误\\
1：格式错误 \\
2：填充错误\\
3：其他错误 \\
(RO)
\end{reglist}\end{regdesc}
\end{register}


\begin{register}{H}{TWAI\_RX\_ERR\_CNT\_REG}{0x0038}\label{regdesc:TWAIRXERRCNTREG}
\regfield{(reserved)}{24}{8}{000000000000000000000000}%
\regfield{TWAI\_RX\_ERR\_CNT}{8}{0}{{0x0}}%
\reglabel{Reset}\regnewline%
\begin{regdesc}\begin{reglist}
\label{fielddesc:TWAIRXERRCNT}\item [TWAI\_RX\_ERR\_CNT] 接收错误计数，数值变化发生在接收状态下。（RO \textcolor{red}{| R/W}）
\end{reglist}\end{regdesc}
\end{register}


\begin{register}{H}{TWAI\_TX\_ERR\_CNT\_REG}{0x003C}\label{regdesc:TWAITXERRCNTREG}
\regfield{(reserved)}{24}{8}{000000000000000000000000}%
\regfield{TWAI\_TX\_ERR\_CNT}{8}{0}{{0x0}}%
\reglabel{Reset}\regnewline%
\begin{regdesc}\begin{reglist}
\label{fielddesc:TWAITXERRCNT}\item [TWAI\_TX\_ERR\_CNT] 发送错误计数，数值变化发生在发送状态下。（RO \textcolor{red}{| R/W}）
\end{reglist}\end{regdesc}
\end{register}


\begin{register}{H}{TWAI\_RX\_MESSAGE\_CNT\_REG}{0x0074}\label{regdesc:TWAIRXMESSAGECNTREG}
\regfield{(reserved)}{25}{7}{0000000000000000000000000}%
\regfield{TWAI\_RX\_MESSAGE\_COUNTER}{7}{0}{{0x0}}%
\reglabel{Reset}\regnewline%
\begin{regdesc}\begin{reglist}
\label{fielddesc:TWAIRXMESSAGECOUNTER}\item [TWAI\_RX\_MESSAGE\_COUNTER] 表示接收 FIFO 中数据包的个数。(RO)
\end{reglist}\end{regdesc}
\end{register}


\begin{register}{H}{TWAI\_INT\_ST\_REG}{0x000C}\label{regdesc:TWAIINTSTREG}
\regfield{(reserved)}{23}{9}{00000000000000000000000}%
\regfield{TWAI\_BUS\_STATE\_INT\_ST}{1}{8}{{0}}%
\regfield{TWAI\_BUS\_ERR\_INT\_ST}{1}{7}{{0}}%
\regfield{TWAI\_ARB\_LOST\_INT\_ST}{1}{6}{{0}}%
\regfield{TWAI\_ERR\_PASSIVE\_INT\_ST}{1}{5}{{0}}%
\regfield{(reserved)}{1}{4}{0}%
\regfield{TWAI\_OVERRUN\_INT\_ST}{1}{3}{{0}}%
\regfield{TWAI\_ERR\_WARN\_INT\_ST}{1}{2}{{0}}%
\regfield{TWAI\_TX\_INT\_ST}{1}{1}{{0}}%
\regfield{TWAI\_RX\_INT\_ST}{1}{0}{{0}}%
\reglabel{Reset}\regnewline%
\begin{regdesc}\begin{reglist}
\label{fielddesc:TWAIRXINTST}\item [TWAI\_RX\_INT\_ST] TWAI\_RX\_INT 的屏蔽中断状态。(RO)
\label{fielddesc:TWAITXINTST}\item [TWAI\_TX\_INT\_ST] TWAI\_TX\_INT 的屏蔽中断状态。(RO)
\label{fielddesc:TWAIERRWARNINTST}\item [TWAI\_ERR\_WARN\_INT\_ST] TWAI\_ERR\_WARN\_INT 的屏蔽中断状态。(RO)
\label{fielddesc:TWAIOVERRUNINTST}\item [TWAI\_OVERRUN\_INT\_ST] TWAI\_OVERRUN\_INT 的屏蔽中断状态。(RO)
\label{fielddesc:TWAIERRPASSIVEINTST}\item [TWAI\_ERR\_PASSIVE\_INT\_ST] TWAI\_ERR\_PASSIVE\_INT 的屏蔽中断状态。(RO)
\label{fielddesc:TWAIARBLOSTINTST}\item [TWAI\_ARB\_LOST\_INT\_ST] TWAI\_ARB\_LOST\_INT 的屏蔽中断状态。(RO)
\label{fielddesc:TWAIBUSERRINTST}\item [TWAI\_BUS\_ERR\_INT\_ST] TWAI\_BUS\_ERR\_INT 的屏蔽中断状态。(RO)
\label{fielddesc:TWAIBUSSTATEINTST}\item [TWAI\_BUS\_STATE\_INT\_ST] TWAI\_BUS\_STATE\_INT 的屏蔽中断状态。(RO)
\end{reglist}\end{regdesc}
\end{register}


\begin{register}{H}{TWAI\_INT\_ENA\_REG}{0x0010}\label{regdesc:TWAIINTENAREG}
\regfield{(reserved)}{23}{9}{000000000000000000000000}%
\regfield{TWAI\_BUS\_STATE\_INT\_ENA}{1}{8}{{0}}%
\regfield{TWAI\_BUS\_ERR\_INT\_ENA}{1}{7}{{0}}%
\regfield{TWAI\_ARB\_LOST\_INT\_ENA}{1}{6}{{0}}%
\regfield{TWAI\_ERR\_PASSIVE\_INT\_ENA}{1}{5}{{0}}%
\regfield{(reserved)}{1}{4}{0}%
\regfield{TWAI\_OVERRUN\_INT\_ENA}{1}{3}{{0}}%
\regfield{TWAI\_ERR\_WARN\_INT\_ENA}{1}{2}{{0}}%
\regfield{TWAI\_TX\_INT\_ENA}{1}{1}{{0}}%
\regfield{TWAI\_RX\_INT\_ENA}{1}{0}{{0}}%
\reglabel{Reset}\regnewline%
\begin{regdesc}\begin{reglist}
\label{fielddesc:TWAIRXINTENA}\item [TWAI\_RX\_INT\_ENA] 置 1 使能接收中断。(R/W)
\label{fielddesc:TWAITXINTENA}\item [TWAI\_TX\_INT\_ENA] 置 1 使能发送中断。(R/W)
\label{fielddesc:TWAIERRWARNINTENA}\item [TWAI\_ERR\_WARN\_INT\_ENA] 置 1 使能错误报警中断。(R/W)
\label{fielddesc:TWAIOVERRUNINTENA}\item [TWAI\_OVERRUN\_INT\_ENA] 置 1 使能数据溢出中断。(R/W)
\label{fielddesc:TWAIERRPASSIVEINTENA}\item [TWAI\_ERR\_PASSIVE\_INT\_ENA] 置 1 使能被动错误中断。(R/W)
\label{fielddesc:TWAIARBLOSTINTENA}\item [TWAI\_ARB\_LOST\_INT\_ENA] 置 1 使能仲裁丢失中断。(R/W)
\label{fielddesc:TWAIBUSERRINTENA}\item [TWAI\_BUS\_ERR\_INT\_ENA] 置 1 使能总线错误中断。(R/W)
\label{fielddesc:TWAIBUSSTATEINTENA}\item [TWAI\_BUS\_STATE\_INT\_ENA] 置 1 使能总线状态中断。(R/W)
\end{reglist}\end{regdesc}
\end{register}
